\documentclass[12pt, a4paper, oneside]{ctexart}
\usepackage{amsmath, amsthm, amssymb, bm, xcolor, tikz, framed, graphicx, hyperref, mathrsfs}

\title{Soild State Physics Homework (Week 9-10)}
\author{Tang Yi \\ [学号已脱敏]}
\usepackage[a4paper, left=2cm, right=2cm, top=2cm, bottom=2cm]{geometry}
\linespread{1.3} % 行距更紧凑
\definecolor{shadecolor}{RGB}{241, 241, 255}

\newcounter{problemname}
\newenvironment{problem}
  {\begin{shaded}\stepcounter{problemname}\par\noindent\textbf{题目 \theproblemname. }}
  {\end{shaded}\par}
\newenvironment{solution}
  {\par\noindent\textbf{题目 \theproblemname 的解答. }}
  {\par}
\newenvironment{note}
  {\par\noindent\textbf{题目 \theproblemname 的注记. }}
  {\par}

\begin{document}
\maketitle

\begin{problem}
根据 $k=\pm \dfrac{\pi}{a}$ 状态简并微扰结果，求出与 $E_+,E_-$ 对应的本征态波函数 $\psi_+,\psi_-$。说明它们都代表驻波，并比较两个电子云分布（即 $|\psi|^2$），说明能隙的来源。假设 $V_K = V_K^*$。
\end{problem}

\begin{solution}
在布里渊区边界：
\[
k=\frac{\pi}{a}
\]
自由电子的两个简并态为
\[
\psi_1=e^{ikx},\qquad \psi_2=e^{-ikx}
\]
周期势引起耦合，其微扰矩阵为
\[
\begin{pmatrix}
E_0 & V_K \\
V_K^* & E_0
\end{pmatrix}
\]
由于题设 $V_K=V_K^*$，故本征值为
\[
E_\pm = E_0 \pm |V_K|
\]

对应本征态：
\[
\psi_+ = \frac{1}{\sqrt2}(e^{ikx}+e^{-ikx})
= \sqrt2 \cos kx
\]
\[
\psi_- = \frac{1}{\sqrt2}(e^{ikx}-e^{-ikx})
= \sqrt2 i \sin kx
\]

二者均为驻波，因为不再具有传播相位因子，而是空间固定振荡。

电子云分布：
\[
|\psi_+|^2 \propto \cos^2 kx
\]
\[
|\psi_-|^2 \propto \sin^2 kx
\]

其中：

\begin{itemize}
\item $\psi_+$ 的电子主要集中在离子附近；
\item $\psi_-$ 的电子主要集中在离子之间。
\end{itemize}

由于离子附近势能较低，因此：
\[
E_+ > E_-
\]
两者分裂形成能隙：
\[
\Delta E = 2|V_K|
\]

这就是布里渊区边界处能隙产生的原因。
\end{solution}

\begin{problem}
写出一维近自由电子近似，第 $n$ 个能带（$n=1,2,3$）中，简约波数 $k=\dfrac{\pi}{2a}$ 时的 0 级波函数。
\end{problem}

\begin{solution}
近自由电子近似中，0级波函数仍为平面波：
\[
\psi_k^{(0)}(x)=e^{ikx}
\]

在简约区方案中：
\[
k=\frac{\pi}{2a}
\]

不同能带对应不同倒格矢平移：
\[
k_n = k + G_n
\]
其中
\[
G_n = \frac{2n\pi}{a}
\]

因此：

第一能带：
\[
\psi_1^{(0)}(x)=e^{i\frac{\pi}{2a}x}
\]

第二能带：
\[
\psi_2^{(0)}(x)=e^{-i\frac{3\pi}{2a}x}
\]

第三能带：
\[
\psi_3^{(0)}(x)=e^{i\frac{5\pi}{2a}x}
\]
\end{solution}

\begin{problem}
电子周期场的势能函数为
\[
V(x)=
\begin{cases}
\dfrac12 m\omega^2\left[b^2-(x-na)^2\right],
& na-b\le x\le na+b
\\
0,
& (n-1)a+b\le x\le na-b
\end{cases}
\]
其中 $a=4b$，$\omega$ 为常数。

\begin{enumerate}
\item 试画出此势能曲线，并求其平均值。
\item 用近自由电子近似模型求出晶体的第二个带隙宽度。
\end{enumerate}
\end{problem}

\begin{solution}
(1) 势能在每个格点附近为抛物线，其余区域为零，因此形成周期性抛物势阱。

平均势能：
\[
\bar V=\frac1a\int_{-b}^{b}\frac12 m\omega^2(b^2-x^2)\,dx
\]

计算：
\[
\bar V
=\frac{m\omega^2}{2a}
\left[
b^2x-\frac{x^3}{3}
\right]_{-b}^{b}
\]

得到：
\[
\bar V
=\frac{m\omega^2}{2a}\cdot\frac{4b^3}{3}
\]

由于 $a=4b$：
\[
\boxed{
\bar V=\frac16 m\omega^2 b^2
}
\]

(2) 第二带隙对应：
\[
G=\frac{4\pi}{a}
\]

带隙宽度：
\[
\Delta = 2|V_G|
\]

其中
\[
V_G=\frac1a\int_{-b}^{b}
\frac12 m\omega^2(b^2-x^2)e^{-iGx}dx
\]

代入
\[
G=\frac{4\pi}{a}=\frac{\pi}{b}
\]

计算可得：
\[
V_G=\frac{2m\omega^2b^2}{\pi^3}
\]

因此第二带隙宽度：
\[
\boxed{
\Delta
=
\frac{4m\omega^2b^2}{\pi^3}
}
\]
\end{solution}

\begin{problem}
用紧束缚近似求出面心立方晶格和体心立方晶格 $s$ 态原子能级相对应的能带 $E(\mathbf{k})$ 关系。
\end{problem}

\begin{solution}
紧束缚近似：
\[
E(\mathbf k)=E_s-\sum_n t_n e^{i\mathbf k\cdot \mathbf R_n}
\]
只考虑最近邻跃迁积分 $t$。

(1) 面心立方（FCC）

最近邻共有 12 个：
\[
\mathbf R=
\frac a2(\pm1,\pm1,0),
\frac a2(\pm1,0,\pm1),
\frac a2(0,\pm1,\pm1)
\]

求和得到：
\[
\boxed{
E(\mathbf k)
=
E_s
-4t
\left[
\cos\frac{k_xa}{2}\cos\frac{k_ya}{2}
+\cos\frac{k_ya}{2}\cos\frac{k_za}{2}
+\cos\frac{k_za}{2}\cos\frac{k_xa}{2}
\right]
}
\]

(2) 体心立方（BCC）

最近邻共有 8 个：
\[
\mathbf R=
\frac a2(\pm1,\pm1,\pm1)
\]

因此：
\[
\boxed{
E(\mathbf k)
=
E_s
-8t
\cos\frac{k_xa}{2}
\cos\frac{k_ya}{2}
\cos\frac{k_za}{2}
}
\]
\end{solution}
\end{document}